\chapter{Introduction}
\label{chap:intro}

% =====================================================================
%  Objectif (cf. INSTRUCTION) : présenter le CONTEXTE du stage (place
%  dans le cursus, entreprises, dates) et SUCCINCTEMENT le sujet.
%  Ne pas hypertrophier l'introduction : le détail technique va dans
%  les chapitres « Travaux effectués ».
% =====================================================================

\section{Contexte de l'alternance}
% TODO : place dans le cursus (Master Informatique, année), nature
% (alternance), période globale, et le fait qu'elle s'est déroulée dans
% DEUX entreprises successives : Free Mobile puis Beemo Technologie.
\emph{[À rédiger.]}

\section{Le fil conducteur : la dette technique}
% TODO : amener le thème. Idée directrice : derrière des missions
% techniques variées (framework de test, portage multi-forks,
% rétro-ingénierie de compilation, développement de feature), un même
% enjeu — maîtriser/réduire la dette technique — et un même argument
% business — vitesse de livraison, fiabilité, coût de maintenance.
\emph{[À rédiger : poser la problématique et les enjeux business.]}

\section{Problématique}
% TODO : formuler la question centrale, p.ex. :
% « Comment des actions d'ingénierie ciblées permettent-elles de réduire
%  la dette technique, et en quoi se traduisent-elles en valeur pour
%  l'entreprise ? »
\emph{[À rédiger.]}

\section{Plan du mémoire}
% TODO : annoncer brièvement le déroulé (entreprises, environnement
% technique, veille sur la dette technique, puis les 4 tâches, conclusion).
\emph{[À rédiger.]}
